\subsection{Root multipliers and all outer connectors}\label{proof:Q5-C066}
\paragraph{Q66.1 Original model and exact old interfaces}

Let \(X\) be a nonempty finite probability space \((X,\allowbreak  \mu)\), with
zero-mass atoms deleted. Assume
\fitdisplay{
K(x,y)=-K(y,x),\qquad |K(x,y)|\le1.
}{}
All expectations, inner products, and adjoints use the same original
measure \(\mu\). Put
\fitdisplay{
Jv=(Ev)1,\quad A=J+K,\quad T=J-K=A^*,\quad
 g=K1,\quad p=1+g,\quad q=1-g,
}{}
\fitdisplay{
f_n=A^{n}\one,\quad \bar f_n=T^{n}\one,\quad h_n=T f_{n-1},\quad
 f_0=\bar f_0=1,\qquad
 I_q(X)=E[q^{-1/3}|X|^{8/3}].
}{Q66.1.1}
Since \(K(x,\allowbreak  x)=0\), \(p(x),\allowbreak  q(x)\ge\mu(x)>0\).
All actual path messages are positive, and \(A,\allowbreak  T\) preserve order.

Write
\fitdisplay{
\delta_j=\|K^{j}\one\|_2^2,\quad d=\delta_1,\quad r=Kg,\quad s=Kr,\quad t=Ks,
}{}
\fitdisplay{
\theta=13/32,\quad c=19/32,\quad k=32/19,\quad \tau=d-\delta_2.
}{Q66.1.2}
The original spectral estimate and elementary moment
relations give
\fitdisplay{
\|K\|^2\le\theta,\quad 0\le d\le1/3,\quad
 d^2\le \delta_2\le\theta d,\quad \delta_{j+1}\le\theta \delta_j,\quad
 cd\le\tau\le d-d^2\le2/9.
}{Q66.1.3}


We use the following exact interfaces.

1. The ordinary path multiplier, from the Q3 operator
   estimate T061 and Q4 C115 formula (8.3), after global sign reversal
   and ordinary adjunction:
\fitdisplay{
\|T M_{\bar f_a}\|_{2\to2}^2\le1-d/72,\qquad a\ge2.
}{Q66.1.4}
   Multiplication and propagation are in precisely this order.

2. The C064 positive-kernel transfer. For a positive operator \(B\)
   with row \(a=B\one>0\),
\fitdisplay{
I_q(BY)^3\le\|B\|_{2\to2}^2E|Y|^8
                       \left(E[a^3/\sqrt q]\right)^2.
}{Q66.1.5}
   Indeed \(|BY|^4\le a^3B(|Y|^4)\), followed by Hölder with
   exponents \(3,\allowbreak  3/2\), proves it. In particular, if
   \(\beta(X)=E[X^2/q]\),
\fitdisplay{
\left(E[\bar f_n^3/\sqrt q]\right)^2
    \le\beta(\bar f_n)E\bar f_n^4,\qquad \beta(\bar f_n)\le1-d\quad(n\ge2).
}{Q66.1.6}
   The latter follows from row Cauchy
   \(\beta(TX)\le E[pX^2]\le\beta(X)\),
   \(\bar f_2=q+r\), and \(r^2\le pqd\).

3. The two Q5-C064 ordinary short norms:
\fitdisplay{
\|T\|^2\le(1+d+k\delta_2)^2\le(1+kd)^2,
    \qquad \|T^2\|^2\le(1+d/2)^4.
}{Q66.1.7}
   The second inequality in the first chain uses \(\delta_2\le\theta d\) and
   \(1+k\theta=k\). These are ordinary \(L^2(\mu)\) norms, not
   a generic weighted two-step contraction.

4. C064 already gives (0.2) for \(a=2\), all \(M\ge1\),
   and for every \(a\ge2\) when \(M=1\) or \(M\ge4\).

5. C058 gives, with the same actual terminal \(qf_b\)
   at every inner length,
\fitdisplay{
E[(pT^L(qf_b))^8]\le1-d/12000,\qquad b\ge2,L\ge1.
}{Q66.1.8}
   C049/C050 give \(1-d/500\) when \(b=2,\allowbreak  L\ge1\), and
   \(1-d/50\) when \(b=2,\allowbreak  L\ge2\).

6. C143 supplies its exact zero-step R directions with
   \(E[q(C^-)^4]\le1\), and actual equal pairs
   \(f_\ell h_\ell\le1\), including their complete reversals.

If \(d=0\), then \(g=0\), \(p=q=1\), and all \(f_n=\bar f_n=1\).
The kernel need not vanish. All actual tree messages are then one
by leaf integration. New expressions involving \(\delta_2/d\) or \(\tau/d\)
below are restricted to \(d>0\); the all-real maps at \(d=0\)
follow directly from (Q66.1.4)--(Q66.1.6).

\paragraph{Q66.2 Full path convolution and a sharper geometric-sum estimate}

Define
\fitdisplay{
S_n=\sum_{j=0}^nK^{j}\one,\qquad G_n(z)=\sum_{j=0}^{n-1}z^j,
}{}
\fitdisplay{
\gamma_1=\gamma_2=0,\qquad
 \gamma_j=\sum_{i=1}^{\lfloor(j-1)/2\rfloor}(-1)^{i-1}\delta_i
 \quad(j\ge3).
}{}
Skewness gives \(ES_n=1-\gamma_{n+1}\), and
\(S_n=AS_{n-1}+\gamma_n1\) gives the full finite identity
\fitdisplay{
S_n=f_n+\sum_{j=3}^n\gamma_j f_{n-j},\qquad
 \tau\le\gamma_j\le d\quad(j\ge3).
}{Q66.2.1}
Every summand is nonnegative. This proves statements for every
integer \(n\), without assuming ordinary path means decrease.
For the same finite skew spectral decomposition,
\fitdisplay{
\|G_n(K)\|\le1\quad(n\ge3),\qquad
 \|S_j\|_2\le1\quad(j\ge2),\qquad 0\le S_n\le2\quad(n\ge3).
}{Q66.2.2}
The scalar verification uses
\fitdisplay{
|G_n(i\sqrt x)|^2=
 \begin{cases}
 (1+x^n)/(1+x),&n\text{ odd},\\
 (1-(-1)^{n/2}x^{n/2})^2/(1+x),&n\text{ even},
 \end{cases}
 \quad0\le x\le\theta.
}{}
For \(n=2\bmod4,\allowbreak  \ n\ge6\), the required inequality follows from
\(2x^2+x^5<1\); the other cases are immediate. Finally
\(|KS_{n-1}|\le\|S_{n-1}\|_2\) proves the pointwise bound in
(Q66.2.2). Positivity in (Q66.2.1) gives \(0<f_n\le2\) for \(n\ge3\).

We need a uniform improvement which retains \(\delta_2\):
\fitdisplay{
\boxed{|G_n(i\sqrt x)|^2\le1-\frac7{15}x,
 \quad n\ge5,\quad0\le x\le\theta.}
}{Q66.2.3}
For every \(n\ge5\), its numerator is at most \((1+x^3)^2\).
For odd \(n\) use \(x^n\le2x^3+x^6\); for \(4\mid n\)
use a numerator at most one; the remaining even case has
\(n/2\ge3\). It therefore suffices to check
\fitdisplay{
8-7x-30x^2-15x^5\ge0.
}{Q66.2.4}
The polynomial decreases on \([0,\allowbreak  \theta]\), and its endpoint is
\fitdisplay{
\frac{1311885}{33554432}>0.
}{}
This proves (Q66.2.3), including \(x=0\).
Integration against the spectral measure at \(g\) gives
\fitdisplay{
\|G_n(K)g\|^2\le d-\frac7{15}\delta_2\quad(n\ge5).
}{Q66.2.5}

A second uniform improvement is needed for \(n\ge6\). Put
\(P(z)=3+2z+z^2\). Then, for \(z=i\sqrt x\),
\fitdisplay{
|P(z)|^2=9-2x+x^2\le9,\qquad
 \operatorname{Re}\{(1-z)^{-1}\overline{P(z)}\}
 =\frac{3+x}{1+x}\ge\frac{109}{45}.
}{}
The geometric tail is bounded by \(3\theta^3=6591/32768\).
Since
\fitdisplay{
\frac{109}{45}-\frac{6591}{32768}-\frac{11}{5}
 =\frac{6217}{294912}>0,
}{}
we obtain the all-length statement
\fitdisplay{
\boxed{\operatorname{Re}\{G_n(z)\overline{P(z)}\}\ge11/5,
 \qquad n\ge6.}
}{Q66.2.6}

\paragraph{Q66.3 Every row length at least six: a uniform joint payment}

\paragraph{Q66.3.1 The actual finite mask and its variance}

For \(n\ge6\), retain the four final terms in (Q66.2.1):
\fitdisplay{
0<f_n\le Z:=S_n-\tau(1+p+f_2+f_3).
}{Q66.3.1}
Here
\fitdisplay{
1+p+f_2+f_3=4-d+3g+2r+s.
}{}
With \(\alpha=\gamma_{n+1}\in[\tau,\allowbreak  d]\),
\fitdisplay{
EZ=1-\alpha-\tau(4-3d)
 \le m:=1-(5-3d)\tau,\qquad Z\le U:=2-\tau.
}{Q66.3.2}
Moreover \(m\ge1/9\): the maximum of
\((d-d^2)(5-3d)\) on \([0,\allowbreak  1/3]\) is \(8/9\), since its
derivative \(5-16d+9d^2\) is positive there.

The nonconstant part of \(Z\) is
\(G_n(K)g-\tau P(K)g\). Equations (Q66.2.5)--(Q66.2.6) give
\fitdisplay{
\|G_n(K)g-\tau P(K)g\|^2
 \le d-\frac7{15}\delta_2-\frac{22}{5}\tau d+9\tau^2d
 \le d-\frac7{15}\delta_2-\frac{12}{5}\tau d,
}{}
where \(9\tau^2\le2\tau\). Its mean is
\(-\alpha+2\tau d\), whose absolute value is at least
\(\tau(1-2d)\). Hence
\fitdisplay{
\operatorname{\operatorname{Var}}Z\le
 V(d,\tau):=d-\frac7{15}\delta_2-\frac{12}{5}\tau d
                   -\tau^2(1-2d)^2.
}{Q66.3.3}
Using \(\delta_2=d-\tau\), this is
\fitdisplay{
V=\frac8{15}d+
       \left(\frac7{15}-\frac{12}{5}d\right)\tau
       -(1-2d)^2\tau^2.
}{Q66.3.4}

This variance majorant is positive over the full continuous region
\fitdisplay{
0<d\le1/3,\qquad cd\le\tau\le d-d^2.
}{Q66.3.5}
Indeed \(\delta_2/d\le\theta\), \(\tau\le2/9\), \(\tau\le d\), and
\(d(1-2d)^2\le2/27\) imply
\fitdisplay{
V/d\ge\frac{389}{480}-\frac8{15}-\frac2{27}
       =\frac{877}{4320}>0.
}{Q66.3.6}
The bound \(d(1-2d)^2\le2/27\) follows by differentiating on
\([0,\allowbreak  1/3]\); the only interior maximum occurs at \(d=1/6\).

For \(0\le W\le U\), with mean \(\mu_W\), the exact identity
\fitdisplay{
EW^4=\mu_W^4+
 E[(W-\mu_W)^2(W^2+2\mu_WW+3\mu_W^2)]
}{}
gives
\fitdisplay{
EW^4\le\mu_W^4+(U^2+2\mu_WU+3\mu_W^2)\operatorname{\operatorname{Var}}W.
}{Q66.3.7}
Its right side is increasing in a nonnegative mean majorant.
Thus, putting \(C=U^2+2mU+3m^2\),
\fitdisplay{
Ef_n^4\le EZ^4\le F(d,\tau):=m^4+C\,V(d,\tau).
}{Q66.3.8}

\paragraph{Q66.3.2 A full-region monotonicity proof in the shared moment}

We prove that \(F(d,\allowbreak  \tau)\) decreases with \(\tau\) throughout
(Q66.3.5). This is a statement about an explicit fourth-moment majorant,
not monotonicity of a mean-plus-variance functional.

Put \(B=5-3d\), so \(B\ge4\), \(m=1-B\tau\), \(U=2-\tau\).
Since \(0\le m\le1\), \(0\le U\le2\),
\fitdisplay{
-C_\tau=2(U+m)+2B(U+3m)\ge4C.
}{Q66.3.9}
For a direct verification, after replacing \(B\) by four the difference
from \(4C\) is
\fitdisplay{
10U+26m-4U^2-8mU-12m^2\ge0.
}{}
For fixed \(m\), this is concave in \(U\), so its minimum is on
\(U=0,\allowbreak  2\). At \(U=0\) it is \(26m-12m^2\ge0\).
At \(U=2\) it is \(4+10m-12m^2\), whose endpoint values on
\(0\le m\le1\) are four and two.

Write \(V=v_0+v_1\tau-w\tau^2\), where
\fitdisplay{
v_0=8d/15,\quad v_1=7/15-(12/5)d,\quad w=(1-2d)^2.
}{}
If \(v_1<0\), then \(V_\tau<0\), and (Q66.3.6), (Q66.3.9) immediately
give \(F_\tau<0\). If \(v_1\ge0\), then
\fitdisplay{
V_\tau-4V
 =v_1-4v_0-4v_1\tau-2w\tau(1-2\tau)
 \le\frac{7-68d}{15}.
}{Q66.3.10}
For \(d\ge7/68\), this again proves \(F_\tau<0\).
For \(0\le d\le7/68\), use \(C\le11\),
\(m\ge1-5d\), and (Q66.3.9) to obtain
\fitdisplay{
F_\tau\le
 -16(1-5d)^3+\frac{11}{15}(7-68d)=:q_0(d).
}{}
On this interval
\fitdisplay{
q_0'(d)=240(1-5d)^2-\frac{748}{15}
 \ge\frac{461648}{69360}>0.
}{}
Thus \(q_0(d)\le q_0(7/68)=-16(33/68)^3<0\).
All cases are covered, proving
\fitdisplay{
F(d,\tau)\le F(d,cd)=:F_*(d).
}{Q66.3.11}
The exact same \(\delta_2\) was retained through (Q66.3.3)--(Q66.3.11).
No independent replacement of its variance contribution by zero
was made.

\paragraph{Q66.3.3 The complete finite scalar certificate}

At \(\tau=cd\), the quantities in (Q66.3.8) are
\fitdisplay{
m=1-5cd+3cd^2,\quad U=2-cd,
}{}
\fitdisplay{
V=\frac{389}{480}d-\left(\frac{12}{5}c+c^2\right)d^2
       +4c^2d^3-4c^2d^4,\qquad c=19/32.
}{Q66.3.12}
Expanding \(F_*=m^4+(U^2+2mU+3m^2)V\), its coefficients after
the constant one are as follows:


\begingroup\small
\setlength{\tabcolsep}{3pt}
\setlength{\LTleft}{0pt}\setlength{\LTright}{\fill}
\begin{longtable}{@{}>{\raggedright\arraybackslash}p{0.1200\linewidth}>{\raggedright\arraybackslash}p{0.4000\linewidth}>{\raggedright\arraybackslash}p{0.4000\linewidth}@{}}
\caption{Coefficients of \(d^j\) in the fourth-moment envelope \(F_*(d)\).}\label{tab:parameters:q5-c066:1}\\
\toprule
\(j\) & Exact coefficient & Convenient upper coefficient \\
\midrule\endfirsthead
\toprule
\(j\) & Exact coefficient & Convenient upper coefficient \\
\midrule\endhead
\bottomrule\endfoot
1 & \(-1421/480\); & \(-59/20\). \\
\addlinespace[2pt]
2 & \(207461/15360\); & \(271/20\). \\
\addlinespace[2pt]
3 & \(-26785934/491520\); & \(-1089/20\). \\
\addlinespace[2pt]
4 & \(575254583/5242880\); & \(2196/20\). \\
\addlinespace[2pt]
5 & \(-1837517436/15728640\); & \(-2336/20\). \\
\addlinespace[2pt]
6 & \(303421583/5242880\); & \(1158/20\). \\
\addlinespace[2pt]
7 & \(-48c^4\); & \(0\). \\
\addlinespace[2pt]
8 & \(-27c^4\); & \(0\). \\
\end{longtable}
\endgroup


Each entry in the last column is at least the exact coefficient.
Let \(\mathcal P(d)\) be the degree-six upper polynomial with constant
one and that column's coefficients. It is nonnegative because it
majorizes \(F_*(d)>0\).

Since \(k=32/19<17/10\),
\fitdisplay{
(1-d)(1+kd)^2F_*(d)
 \le (1-d)(1+17d/10)^2\mathcal P(d).
}{Q66.3.13}
The factor on the right outside \(\mathcal P\) is
\(1+(12/5)d-(51/100)d^2-(289/100)d^3\).
The right side of (Q66.3.13), minus one, is \(dQ(d)/2000\), where
\fitdisplay{
\begin{aligned}
 Q(d)={}&-1100+11920d-46631d^2-38530d^3+270660d^4\\
       &-242115d^5-237588d^6+616046d^7-334662d^8.
 \end{aligned}
}{Q66.3.14}
For \(d\le1/3\),
\fitdisplay{
-237588d^6+616046d^7\le-\frac{96718}{3}d^6\le0.
}{}
Dropping this pair and the final negative term leaves the
degree-five upper polynomial
\fitdisplay{
Q_5=-1100+11920d-46631d^2-38530d^3+270660d^4-242115d^5.
}{}
For \(z=3d\in[0,\allowbreak  1]\), its degree-five Bernstein coefficients are
\fitdisplay{
-1100,\quad-\frac{916}{3},\quad-\frac{2591}{90},\quad
 -\frac{111529}{270},\quad-\frac{755394}{810},\quad
 -\frac{112574}{81}.
}{Q66.3.15}
Each is at most \(-2591/90\).
Here a Bernstein list \((B_i)_{i=0}^n\) means the exact identity
\fitdisplay{
Q_5(z/3)=\sum_{i=0}^nB_i\binom ni z^i(1-z)^{n-i}.
}{}
Consequently, with \(\eta=2591/180000\),
\fitdisplay{
\boxed{
 (1-d)(1+kd)^2Ef_n^4\le1-\eta d,\qquad n\ge6.}
}{Q66.3.16}
Global sign reversal gives the same bound for \(E\bar f_n^4\).
The \(d=0\) boundary is direct. These are common bounds under sign
reversal, not an equality of opposite fourth moments for a fixed
kernel.

\paragraph{Q66.4 The remaining row length five}

We need a stronger global fifth-path fourth budget than the one
used in C064:
\fitdisplay{
\boxed{Ef_5^4,\ E\bar f_5^4\le1-\frac32d.}
}{Q66.4.1}
The exact convolution gives
\fitdisplay{
f_5=S_5-d(f_2+p+1)+\delta_2
     =S_5-d(f_2+p)-\tau.
}{}
Therefore
\fitdisplay{
Ef_5=1-4d+2\delta_2+d^2\le m:=1-\frac{51}{16}d+d^2,\qquad
 0<f_5\le U:=2-cd.
}{Q66.4.2}
The mean majorant is positive throughout the interval; at \(d=1/3\)
it is \(7/144\), and it is decreasing on the whole interval.

The nonconstant polynomial is \(G_5(K)g-d(2g+r)\).
For \(z=i\sqrt x\),
\fitdisplay{
|G_5(z)|^2\le1-\frac7{15}x,\quad
 \operatorname{Re}\{G_5(z)\overline{(2+z)}\}
 =2-x+x^2\ge7/4,\quad |2+z|^2\le9/2.
}{}
Thus, dropping the nonnegative square of the mean,
\fitdisplay{
\operatorname{\operatorname{Var}}f_5
 \le d-\frac7{15}\delta_2-\frac72d^2+\frac92d^3
 \le d-\frac{37}{15}d^2=:V_5.
}{Q66.4.3}
The last quantity is nonnegative, and positive for \(d>0\).
Apply (Q66.3.7) with (Q66.4.2)--(Q66.4.3).
The coefficients of the resulting polynomial after its constant one
are


\begingroup\small
\setlength{\tabcolsep}{3pt}
\setlength{\LTleft}{0pt}\setlength{\LTright}{\fill}
\begin{longtable}{@{}>{\raggedright\arraybackslash}p{0.1200\linewidth}>{\raggedright\arraybackslash}p{0.4000\linewidth}>{\raggedright\arraybackslash}p{0.4000\linewidth}@{}}
\caption{Coefficients of \(d^j\) in the fifth-path fourth-moment envelope.}\label{tab:parameters:q5-c066:2}\\
\toprule
\(j\) & Exact coefficient & Convenient upper coefficient \\
\midrule\endfirsthead
\toprule
\(j\) & Exact coefficient & Convenient upper coefficient \\
\midrule\endhead
\bottomrule\endfoot
1 & \(-7/4\); & \(-7/4\). \\
\addlinespace[2pt]
2 & \(4589/1920\); & \(12/5\). \\
\addlinespace[2pt]
3 & \(-91549/2560\); & \(-143/4\). \\
\addlinespace[2pt]
4 & \(99070783/983040\); & \(504/5\). \\
\addlinespace[2pt]
5 & \(-352321/3072\); & \(-573/5\). \\
\addlinespace[2pt]
6 & \(36839/640\); & \(288/5\). \\
\addlinespace[2pt]
7 & \(-51/4\); & retained with degree eight. \\
\addlinespace[2pt]
8 & \(1\); & retained with degree seven. \\
\end{longtable}
\endgroup


The final two terms sum to \(d^7(-51/4+d)\le0\).
After discarding them, the difference from \(1-3d/2\), divided
by \(d>0\), is at most
\fitdisplay{
Q_{\rm five}(d)
 =-\frac14+\frac{12}{5}d-\frac{143}{4}d^2
       +\frac{504}{5}d^3-\frac{573}{5}d^4+\frac{288}{5}d^5.
}{Q66.4.4}
The degree-five Bernstein coefficients of \(Q_{\rm five}(z/3)\)
are
\fitdisplay{
-\frac14,\quad-\frac9{100},\quad-\frac{589}{1800},\quad
 -\frac{353}{600},\quad-\frac{1057}{1350},\quad-\frac{13}{15}.
}{Q66.4.5}
They are all negative. This proves (Q66.4.1) on the whole domain,
including \(d=0\) separately. The improved pointwise upper bound
\(f_5\le2-\tau\), which follows from the exact positive convolution,
is retained in this proof.

\paragraph{Q66.5 Additional strict ordinary loss for the actual root \(\bar f_3\)}

We prove the strengthened ordinary multiplier
\fitdisplay{
\boxed{
 \|T M_{\bar f_3}\|^2\le(1-d/2)^2(1-d/72).}
}{Q66.5.1}
We use the \(S_3\) case of the Q3 T061 Schur argument below.

Set
\fitdisplay{
H=S_3=1+g+r+s,\quad a=H-1,\quad
 R=(I+K^2)^{-1},\quad e=-R(K^4+K^5)1,\quad u=1+e.
}{}
The inverse exists by \(\|K\|^2\le13/32<1\).
Let
\fitdisplay{
P=T A,\quad t_0=1-EH=d,\quad v=E(gH),\quad
 b=E(He),\quad c_0=E(ge),\quad a_0=E|ae|.
}{}
The elementary exact formulas are
\fitdisplay{
Pz=qEz-E(gz)-K^2z,\qquad
 (I-K^2)a-g-r=(I+K^2)e.
}{Q66.5.2}
Expanding by (Q66.5.2) gives
\fitdisplay{
M_H P M_Hu-u
 =-a^2+ae+H\{q(b-t_0)-(v+c_0)+T K(ae)\}.
}{Q66.5.3}
For a direct check, write
\(E(Hu)=1-t_0+b\),
\(E(gHu)=v+c_0+E(gae)\), and
\(K^2(Hu)=K^21+K^2a+K^2e+K^2(ae)\).
The last two \(ae\) contributions form \(TK(ae)\);
the remaining terms inside \(P(Hu)\) reduce to \(1-a+e\).
Multiplying by \(H=1+a\) and subtracting \(u\) yields (Q66.5.3).

Let \(\nu\) be the finite spectral measure of \(-K^2\) at \(1\).
Using skewness and the explicit powers in \(H,\allowbreak  e\),
\fitdisplay{
b=-\delta_2-\delta_3\le0,\qquad
 v+c_0=\int\frac{x(1-2x)}{1-x}\,d\nu(x)\ge\frac27d.
}{Q66.5.4}
For the second bound use \(x\le\theta<\rho:=5/12\) and
\((1-2\rho)/(1-\rho)=2/7\).
Also \(\|a\|_2\le\sqrt d\), since
\(a=G_3(K)g\) and \(|G_3(i\sqrt x)|^2=(1+x^3)/(1+x)\le1\).

The exact spectral formulas for \(e\) give
\fitdisplay{
\|e\|_2^2=\int\frac{x^4(1+x)}{(1-x)^2}\,d\nu(x)
 \le\frac{2125}{7056}d\le\frac{16}{49}d,
}{}
and, writing \(e=-KR(K^3+K^4)1\),
\fitdisplay{
|e|^2\le\|R(K^3+K^4)1\|_2^2
 \le\frac{425}{588}d\le\frac{36}{49}d.
}{Q66.5.5}
The two rational bounds are obtained by evaluating the increasing
scalar factors at \(\rho=5/12\). They imply
\fitdisplay{
a_0\le4d/7,\quad |e|\le6\sqrt d/7\le18/35,\quad
 17/35\le u\le53/35.
}{Q66.5.6}
Here \(\sqrt d\le3/5\) follows from \(d\le1/3\).

Pointwise \(|K(ae)|\le a_0\), and positivity of \(T\) gives
\(TK(ae)\le q a_0\). Since \(b-d+a_0\le0\), (Q66.5.3)--(Q66.5.4) imply
\fitdisplay{
M_H P M_Hu-u
 \le-a^2+a(e-2d/7)-2d/7.
}{}
Moreover \(|e-2d/7|\le(36/35)\sqrt d\). Completing the square,
\fitdisplay{
M_H P M_Hu-u
 \le\frac{324}{1225}d-\frac27d
 =-\frac{26}{1225}d.
}{Q66.5.7}
Because \(u\le53/35\), this is at most
\(-26du/1855\).

The operator \(M_H P M_H\) is positive semidefinite and has a
nonnegative kernel. Indeed \(P=T A=A^*A\), and \(H=f_3+d>0\).
The elementary weighted Schur argument therefore gives
\fitdisplay{
\|M_H T\|^2=\|M_H P M_H\|\le1-\frac{26}{1855}d.
}{Q66.5.8}
For clarity, if a symmetric nonnegative kernel \(Q\) satisfies
\(Qu\le\lambda u\) with \(u>0\), the inequality
\(2z(x)z(y)\le(u(y)/u(x))z(x)^2+(u(x)/u(y))z(y)^2\), integrated against
\(Q(x,\allowbreak  y)\mu(x)\mu(y)\), gives
\(\langle z,\allowbreak  Qz\rangle\le\lambda\|z\|^2\).
Positive semidefiniteness then controls the full operator norm.

Finally \(f_3=H-d\), \(0<H\le2\), so
\fitdisplay{
0<f_3/H\le1-d/2.
}{}
Thus
\fitdisplay{
\|M_{f_3}T\|^2\le(1-d/2)^2(1-26d/1855).
}{}
Global sign reversal followed by the ordinary adjoint gives the
same estimate for \(T M_{\bar f_3}\). Since \(26/1855\ge1/72\),
(Q66.5.1) follows. 

\paragraph{Q66.6 All short outer lengths for every \(a\ge3\)}

Put \(B_{M,\allowbreak  a}=T^M M_{\bar f_a}\); its actual row is \(\bar f_{M+a}\).
Using (Q66.1.4)--(Q66.1.6),
\fitdisplay{
I_q(B_{M,a}Y)^3
 \le(1-d/72)\|T^{M-1}\|^2(1-d)E\bar f_{M+a}^4\,E|Y|^8.
}{Q66.6.1}

If \(M=2,\allowbreak  a\ge4\), the row length is at least six. Equations
(Q66.1.7) and (Q66.3.16) show
\fitdisplay{
I_q(T^2(\bar f_aY))^3
 \le(1-d/72)(1-\eta d)E|Y|^8
 \le(1-d/72)E|Y|^8.
}{Q66.6.2}
If \(M=3,\allowbreak  a\ge3\), again the row length is at least six, and
\fitdisplay{
(1+d/2)^2\le1+\frac{13}{12}d\le1+kd.
}{}
Therefore the same payment proves
\fitdisplay{
I_q(T^3(\bar f_aY))^3
 \le(1-d/72)(1-\eta d)E|Y|^8.
}{Q66.6.3}
This covers every such \(a\), not merely sufficiently long roots.

It remains to consider \(a=3,\allowbreak  M=2\). Use (Q66.5.1) in place of the
weaker ordinary multiplier in (Q66.6.1), and use (Q66.4.1):
\fitdisplay{
I_q(T^2(\bar f_3Y))^3
 \le(1-d/72)(1-d/2)^2(1+kd)^2(1-d)(1-3d/2)E|Y|^8.
}{Q66.6.4}
All displayed factors are nonnegative. We have
\fitdisplay{
(1-d)(1-3d/2)\le(1-5d/4)^2,
}{}
and \(k<7/4\), so the square root of the extra factor in (Q66.6.4)
is at most
\fitdisplay{
(1-d/2)(1+7d/4)(1-5d/4)
 =1-\frac{39}{16}d^2+\frac{35}{32}d^3\le1.
}{Q66.6.5}
The last step uses \(d\le1/3\).
This proves the final short case on the whole continuous domain.

The regular case follows directly from (Q66.1.4)--(Q66.1.6).
Equations (Q66.6.2)--(Q66.6.5) prove (0.1). Together with the exact
C064 cases, they prove the complete all-\(a\), all-\(M\) map (0.2).
The argument uses the operator order
\(T^{M-1}(T M_{\bar f_a})\), the row \(\bar f_{M+a}\), and the
fixed measure.


\paragraph{Q66.10.1 A complete outer-length interface for every negative even root}

Let an actual even residual root table be
\fitdisplay{
1\le s_1<t_1<s_2<t_2<\cdots<s_j<t_j,\qquad j\ge1.
}{}
The negative canonical pair orientation has factor
\fitdisplay{
C_\Sigma^-=E_r\prod_{i=1}^j\bar h_{s_i}\bar f_{t_i},
 \qquad \bar h_s=A \bar f_{s-1}.
}{Q66.10.1}
Thus each shorter arm has its first edge away from the root and all
remaining edges toward the root; each longer arm is wholly toward
the root. Equal pairs in \(E_r\) are the actual reversed pairs.
Every original arm and its length is retained.

Choose the first longer arm \(t_1\ge2\) and set
\fitdisplay{
W_\Sigma=E_r\bar h_{s_1}
                     \prod_{i=2}^j\bar h_{s_i}\bar f_{t_i}.
}{}
For any fixed actual terminal message \(F\) and any \(L,\allowbreak  M\ge1\),
the actual half-tree factor and output are
\fitdisplay{
H_\Sigma=C_\Sigma^-T^L F
          =\bar f_{t_1}(W_\Sigma T^L F),\qquad T^M H_\Sigma.
}{}
The full multiplier theorem proves the unconditional analytic bound
\fitdisplay{
\boxed{
 I_q(T^M H_\Sigma)^3
 \le(1-d/72)E[(W_\Sigma T^L F)^8],\qquad M\ge1.}
}{Q66.10.2}
There is no short-outer exception left in this interface, and
\(L,\allowbreak  M\) are independent. The missing responsibility for these more
general roots is now an explicit ordinary eighth input, not an
unspecified outer-propagation claim.

For the terminal family
\(F=E_tqf_b\), \(b\ge2\), the precise sufficient new input is
\fitdisplay{
E\!\left[\left\{
 E_r(A \bar f_{s_1-1})
 \left(\prod_{i=2}^j(A \bar f_{s_i-1})\bar f_{t_i}\right)
 T^L(E_tqf_b)\right\}^{8}\right]\le1.
}{Q66.10.3}
When \(j=1,\allowbreak  s_1=1\), this is exactly the input with factor
\(p=A1\), and parts Q66.8--9 finish all outer lengths. For \(s_1\ge2\), or \(j\ge2\), (Q66.10.3) is a sufficient local condition, not a necessary condition for
the full mapping or for the existence of an orientation of the
original tree.

The following two symbolic examinations show why the missing factors matter in the current positive comparison.
